\emph{Sistem bukti} mènèhi cara sintaktis kanggo mènèhi ciri
konsèkuènsi lan kompatibilitas antarané ukara. \emph{Kalkulus
sekèn} iku salah sijiné. \emph{Derivasi} ing kalkulus iki
awujud wit sekèn. Sekèn $\Gamma\Sequent\Delta$ dumadi saka
rong urutan formula kang dipisah~$\Sequent$. Sekèn paling
ndhuwur ing !!a{derivation} iku \emph{sekèn wiwitan} kang
awujud $!A\Sequent !A$. Supaya !!{derivation} bener, sekèn
liyané kudu didhasaraké marang salah siji \emph{paugeran
inferensi}. Paugeran iku pepasangan: siji kanggo sisih kiwa lan
siji kanggo sisih tengen sekèn, kanggo saben panggandhèng lan
kuantor. Upamané, yèn $\Gamma\Sequent\Delta,!A\lif !B$ dijupuk
nganggo paugeran $\RightR{\lif}$, sekèn sadurungé, yaiku
\emph{prèmisé}, kudu $!A,\Gamma\Sequent\Delta,!B$.
Sawetara paugeran uga bisa ngowahi urutan utawa cacah formula
ing sekèn. Paugeran \RightR{\Exchange} ngidinaké rong formula
ing sisih tengen sekèn diijolaké.

Yèn ana !!a{derivation} kanggo sekèn $\quad\Sequent !A$,
$!A$ diarani \emph{teorema}, lan kita nulis $\Proves !A$.
Yèn ana !!a{derivation} saka $\Gamma_0\Sequent !A$ kanthi
saben $!B$ ing $\Gamma_0$ uga ana ing~$\Gamma$, $!A$
\emph{bisa diderivasi saka}~$\Gamma$, ditulis $\Gamma\Proves !A$.
Yèn ana !!a{derivation} saka $\Gamma_0\Sequent\quad$ kanthi
saben $!B$ ing $\Gamma_0$ uga ana ing~$\Gamma$, $\Gamma$
\emph{ora konsistèn}. Yèn ora, $\Gamma$ \emph{konsistèn}.
Konsèp-konsèp iki gegayutan. Upamané, $\Gamma\Proves !A$
yèn lan mung yèn $\Gamma\cup\{\lnot !A\}$ ora konsistèn.
Uga ana gegayutan karo semantik. Yèn $\Gamma\Proves !A$,
mula $\Gamma\Entails !A$. Sipat yèn apa kang bisa diderivasi
saka~$\Gamma$ mesthi dadi konsèkuènsi saka~$\Gamma$ diarani
\emph{kesahihan}. \emph{Teorema kesahihan} dibuktèkaké kanthi
induksi tumrap dawané !!{derivation}s. Saben inferensi njaga
kesahan sekèn kesimpulan yèn sekèn prèmis sah.
